% Generated draft: requires mathematical and source review.
\documentclass[11pt]{article}
\usepackage[margin=25mm]{geometry}
\usepackage{fontspec}
\setmainfont{Latin Modern Roman}
\usepackage{amsmath,amssymb,mathtools,unicode-math}
\setmathfont{Latin Modern Math}
\usepackage{enumitem,needspace,xurl,hyperref}
\setlength{\parindent}{0pt}
\setlength{\parskip}{5pt}
\setlength{\emergencystretch}{3em}
\newcommand{\N}{\mathbb{N}}
\newcommand{\Z}{\mathbb{Z}}
\newcommand{\Q}{\mathbb{Q}}
\newcommand{\R}{\mathbb{R}}
\newcommand{\C}{\mathbb{C}}
\newcommand{\F}{\mathbb{F}}
\newcommand{\A}{\mathbb{A}}
\newcommand{\PP}{\mathbb{P}}
\newcommand{\E}{\mathbb{E}}
\newcommand{\eps}{\varepsilon}
\newcommand{\dd}{\,\mathrm{d}}
\begin{document}
\section*{UnsolvedMath ID: 20000534}
\textbf{Source problem code:} catalogue entry
\subsection{Definitions and mathematical statement}
Let \(\sigma\) be probability arclength measure on \(S^1\subset\R^2\), and let \(R_\theta\) be rotation through angle \(\theta\). For a nonempty radius set \(E\subset(0,\infty)\), define, initially for smooth compactly supported functions,
\[
\mathcal M_{E,\theta}(f,g)(x)=\sup_{t\in E}
\left|\int_{S^1}f(x-tw)g(x-tR_\theta w)\,\dd\sigma(w)\right|.
\]
Krause's Question 7 concerns \(E=(0,\infty)\) and \(\theta=\pi/3\): find exponent triples for which there is a finite constant independent of \(f,g\) such that
\[
\|\mathcal M_{(0,\infty),\pi/3}(f,g)\|_{L^r(\R^2)}
\leq C_{p,q,r}\|f\|_{L^p(\R^2)}\|g\|_{L^q(\R^2)}.
\]
The source says only ``for certain \(p,q,r\),'' without an endpoint or full-range conjecture. Finite positive exponents use their usual \(L^p\) norms or quasi-norms, with the essential-supremum convention at \(\infty\). For any scale-invariant continuous-radius bound, the natural scaling relation is \(1/r=1/p+1/q\), with \(1/\infty=0\).

The workshop text asserts that this question is equivalent to the corresponding estimate for
\[
\mathcal M_{(0,\infty),\pi}(f,g)(x)=\sup_{t>0}
\left|\int_{S^1}f(x-tw)g(x+tw)\,\dd\sigma(w)\right|.
\]
The equivalence is retained as a source assertion; no transformation, constants, or precise endpoint scope is supplied. It is not an asserted identity between the two operators.

The attached aid describes a linear-maximal transfer. Put
\[
\mathcal C_E h(x)=\sup_{t\in E}\int_{S^1}|h(x-tw)|\,\dd\sigma(w).
\]
For conjugate exponents \(a,a'>1\), \(1/a+1/a'=1\), it records the pointwise estimate
\[
\mathcal M_{E,\theta}(f,g)
\leq[\mathcal C_E(|f|^a)]^{1/a}
[\mathcal C_E(|g|^{a'})]^{1/a'}.
\]
Using the appropriate linear circular maximal bounds, the aid reports the known continuous-radius estimates on the H\"older line with \(r>2\). For the dyadic radius set \(E=\{2^j:j\in\Z\}\), it reports strong bounds for all
\[
1<p,q\leq\infty,\quad1<r\leq\infty,\qquad
\frac1r=\frac1p+\frac1q<1,
\]
uniformly in every rotation \(\theta\), including \(\pi/3\) and \(\pi\). The dyadic result does not replace the continuous supremum in the original problem. The aid identifies the continuous-radius region \(p,q>2\), \(1<r\leq2\), on the same H\"older line as unresolved in the literature considered. It makes no claim of novelty for the elementary pointwise H\"older inequality itself. These are supplied unreviewed partial-result and status claims, not a new sharp-range theorem.
\subsection{Short English statement}
Find \(L^p\times L^q\to L^r\) estimates for a continuous-radius bilinear circular maximal operator with a sixty-degree rotation. The source also asserts an antipodal equivalence. The attached aid gives a transfer from linear circular maximal bounds and a larger dyadic-radius range, which leaves the stated continuous-radius question distinct.
\subsection{Sources}
\begin{itemize}
\item Krause, Question 7, equations (0.7) and (0.8), AIM workshop \emph{Carleson theorems and multilinear operators}: \url{https://aimath.org/pastworkshops/multilinopsproblems.pdf}. The supplied verification says the damaged displays were checked against the official AIM TeX source as well as the PDF; no separate TeX URL is supplied.
\item Local registry \texttt{data/unsolvedmath/problems.json}, permanent ID \(20000534\), exact source code \texttt{catalogue entry}; canonical \texttt{aim-analysis-notes.json}, notes[65].
\item The attached research attempt supplies the radius-set transfer and dyadic result but no full external references for the linear maximal inputs. Attempt 1 is explicitly machine-generated, not peer reviewed, and requires independent verification before citation.
\end{itemize}
\end{document}
